\section*{Chapitre \ref{ch:rs:capacity-decay-and-crowding} --- Capacité, décroissance et encombrement}

\begin{solution}{exo:rs:capacity-decay-and-crowding:1}
$k\sqrt{10^{10}} = 0.10$ donne $k = 10^{-6}$~; $A^* = (2 \times 0.10/(3 \times 10^{-6}))^2 = 4.44$~milliards de dollars, où le rendement net vaut $0.10/3 = 3.3~\%$.
\end{solution}

\begin{solution}{exo:rs:capacity-decay-and-crowding:2}
Le rendement du capital baisse dès le premier dollar, mais le profit en dollars est égal à la taille multipliée par le rendement~: il continue de croître tant que le rendement baisse moins vite que la taille
croît, jusqu'au point où le dollar marginal ne rapporte rien ($g - \frac32 k\sqrt A = 0$).
\end{solution}

\begin{solution}{exo:rs:capacity-decay-and-crowding:3}
Une position de 2~\% de 2~milliards de dollars, soit 40~millions de dollars, vendue sur trois jours~: 13.3~millions de dollars par jour contre 200~millions de dollars échangés, représentant 6.7~\% du volume. L'impact est
$0.7 \times 2~\% \times \sqrt{0.067} = 0.36~\%$ par jour.
\end{solution}

\begin{solution}{exo:rs:capacity-decay-and-crowding:4}
L'impact sur chaque titre est proportionnel à la racine carrée de la transaction du vendeur sur ce titre, et celle-ci est proportionnelle à la fraction vendue~: la perte d'un fonds détenant
le même titre varie comme sa racine carrée. La perte d'un fonds qui ne vend pas est sa position multipliée par l'impact, donc proportionnelle à la part de sa position située là où le vendeur
opère~: il s'agit du recouvrement.
\end{solution}

\begin{solution}{exo:rs:capacity-decay-and-crowding:5}
La perte s'est limitée aux portefeuilles long--short sur facteurs, et non au marché~: quelqu'un vendait la jambe longue et rachetait la jambe courte des portefeuilles valeur et momentum
au même moment. C'est la signature d'un désendettement de fonds détenant des portefeuilles semblables, comme l'ont conclu Khandani et Lo, et non d'une information nouvelle.
\end{solution}

\begin{solution}{exo:rs:capacity-decay-and-crowding:6}
Les coûts croissent plus vite que la taille~: le préjudice de la taille vient donc des transactions. Un portefeuille réoptimisé à chaque taille effectue moins de transactions à mesure qu'il grossit (à 10~milliards de dollars, il renouvelle 1.8~\% de
son capital par jour, contre 227~\% à 10~millions de dollars au chapitre~27), ce qui empêche ses coûts de dépasser son rendement~; le signal était le même.
\end{solution}

\begin{solution}{exo:rs:capacity-decay-and-crowding:7}
Avec des portefeuilles identiques, le quinzième jour après le pic vaut $1.30/3.07 = 0.42$ pour l'ensemble du portefeuille et $0.41/0.97 = 0.42$ pour un dixième~: au-dessus de la part permanente de 0.3
parce que le pic, au troisième jour, intègre déjà une partie de la décroissance de l'\omterm{def:rs:transaction-costs-in-research:impact}{impact temporaire} des premiers jours.
\end{solution}

\begin{solution}{exo:rs:capacity-decay-and-crowding:8}
Un ratio de Sharpe positif n'est pas une \omterm{def:rs:capacity-decay-and-crowding:capacity}{capacité}~: à 5~milliards de dollars, le ratio de la règle fixe est tombé au cinquième de son meilleur niveau (0.34 contre 1.79) et son profit en dollars est
passé son maximum (146~millions de dollars à 3~milliards, 122~millions de dollars à 5~milliards). Selon le critère du chapitre, sa \omterm{def:rs:capacity-decay-and-crowding:capacity}{capacité} est de 2.02~milliards de dollars, et l'estimation est elle-même
incertaine~; il faut la faire tourner en dessous de ce niveau.
\end{solution}

\begin{solution}{pb:rs:capacity-decay-and-crowding:1}
\begin{enumerate}
\item $A^* = (2g/3k)^2$, où $g - \frac32 k\sqrt A = 0$~; le rendement net y vaut $g/3$.
\item 1.79, 1.74, 1.64, 1.48, 1.18, 0.90, 0.68, 0.34, $-0.28$, $-1.85$ à 10~M\$, 30~M\$, 100~M\$, 300~M\$, 1~Md\$, 2~Md\$, 3~Md\$, 5~Md\$, 10~Md\$, 30~Md\$.
\item Le profit atteint son maximum à 3~milliards de dollars (146~millions de dollars par an)~; le ratio de Sharpe est divisé par deux à 2.02~milliards de dollars.
\item Le portefeuille effectue moins de transactions à mesure qu'il grossit~: ratio de Sharpe de 2.43 à 1~milliard de dollars, 1.36 à 30~milliards de dollars, le profit continuant de croître (806~millions de dollars).
\item En deçà de la \omterm{def:rs:capacity-decay-and-crowding:capacity}{taille maximisant le profit}, là où la courbe est plate, car la courbe va bouger (décroissance, concurrents) et l'estimation est incertaine.
\item Le signal se dégrade à mesure que d'autres le découvrent, et les coûts augmentent à mesure que d'autres négocient les mêmes titres.
\item L'\omterm{def:rs:capacity-decay-and-crowding:crowding}{encombrement} (crowding)~: un arbitragiste ignore combien de pairs se trouvent dans la même position, et l'effet de levier, choix privé qui crée une externalité de vente forcée.
\item La corrélation moyenne des rendements anormaux entre les actions d'une stratégie momentum~; des valeurs élevées annonçaient de fortes inversions des rendements momentum.
\item La corrélation entre les titres propres à la stratégie, nette des facteurs, le recouvrement avec les portefeuilles estimés d'autres acteurs, l'encours vendeur à découvert et les jours de couverture, ainsi que la liquidité du portefeuille.
\item Cinq fonds de 2~milliards de dollars, exposition brute égale à quatre fois le capital, 100 titres de chaque côté parmi 500 actions échangeant 200~millions de dollars par jour avec une volatilité de 2~\%~; l'un d'eux vend sur trois jours~;
impact en racine carrée avec $\eta = 0.7$, 30~\% permanent, le reste se résorbant avec une demi-vie d'un jour.
\item De $-0.03~\%$ (10~\% vendus, recouvrement 0.04) à $-3.07~\%$ (tout vendu, portefeuilles identiques)~; le tableau du texte.
\item Au troisième jour, dernier jour de vente~; au quinzième jour, le cas de l'ensemble du portefeuille est remonté à $-1.30~\%$.
\item Du 6 au 9 août, la valeur a reculé de 2.28~\% et le momentum de 3.43~\% (10.6 et 12.5 écarts-types quotidiens), tandis que le marché progressait de 1.39~\%~; le 10 août, la valeur a regagné
1.42~\% et le momentum de 1.15~\%.
\item Pertes de portefeuilles long--short de type valeur, débutées en juillet, deux \omterm{def:rs:capacity-decay-and-crowding:crowding}{débouclages} les 1 et 6 août, commençant par les financières, acheteurs sur le book-to-market et vendeurs sur le momentum,
et un retrait du capital des teneurs de marché à partir du 8 août.
\item \textbf{Résultat central.} Avec des portefeuilles identiques, la perte maximale des fonds non vendeurs vaut 0.97~\%, 1.54~\%, 2.17~\% et 3.07~\% du capital lorsque sont vendus respectivement un dixième, un quart, la moitié et
la totalité du portefeuille (proportionnellement moins à faible recouvrement)~; le ratio de Sharpe net de la règle fixe est divisé par deux à 2.02~milliards de dollars.
\item Ne pas vendre au plus bas s'il peut tenir~: la part temporaire revient~; ne réduire que dans la mesure où ses propres limites l'exigent, et fournir de la liquidité s'il dispose de capital.
\item Elles transforment les pertes en ventes supplémentaires~: les pics s'approfondissent et gagnent les fonds qui ne recouvraient pas le premier vendeur.
\item Les transactions de ses stratégies sur les mêmes titres s'additionnent sous une loi de coût concave, si bien que leur \omterm{def:rs:capacity-decay-and-crowding:capacity}{capacité} combinée est inférieure à la somme.
\item La \omterm{def:rs:capacity-decay-and-crowding:capacity}{courbe de capacité}, selon les coûts de la société, la rotation et la durée de détention, la liquidité des titres, et les autres acteurs susceptibles de détenir le même portefeuille.
\item Ses coûts à mesure qu'il grossit, déterminés par les volumes et les titres qu'il échange, et par le capital des autres acteurs engagés dans la même position.
\end{enumerate}
\end{solution}

\begin{solution}{iq:rs:capacity-decay-and-crowding:1}
Pendant la semaine du 6 août 2007, les portefeuilles quantitatifs long--short sur actions ont fortement reculé sans nouvelle particulière, alors que les facteurs valeur et momentum chutaient de dix écarts-types quotidiens ou plus
tandis que le marché progressait~; le 10 août, une grande partie de ce mouvement s'est inversée. Khandani et Lo l'attribuent au \omterm{def:rs:capacity-decay-and-crowding:crowding}{débouclage} de portefeuilles semblables par plusieurs fonds et à un
retrait du capital des teneurs de marché.
\end{solution}

\begin{solution}{iq:rs:capacity-decay-and-crowding:2}
Relancer le \omterm{def:rs:vectorised-backtests:backtest}{backtest} sur une grille de tailles, avec le modèle de coûts estimé par la société (fourchette et impact en racine carrée), réoptimisé à chaque taille~; lire le ratio de Sharpe net et
le profit en dollars selon la taille~; retenir la \omterm{def:rs:capacity-decay-and-crowding:capacity}{taille maximisant le profit} et un critère, tel que la taille à laquelle le ratio de Sharpe net est divisé par deux~; refaire le calcul avec des données en temps réel.
\end{solution}

\begin{solution}{iq:rs:capacity-decay-and-crowding:3}
Rechercher une corrélation excédentaire entre les titres de la stratégie, nette des facteurs du modèle de risque (\omterm{def:rs:capacity-decay-and-crowding:crowding}{comomentum}), le recouvrement avec les portefeuilles estimés d'autres acteurs (déclarations de détention, rendements
de produits semblables), un encours vendeur à découvert et des jours de couverture élevés du côté vendeur, et vérifier si les pertes de la stratégie coïncident avec celles des autres.
\end{solution}

\begin{solution}{iq:rs:capacity-decay-and-crowding:4}
Estimer le recouvrement et la taille du vendeur~; prévoir une perte temporaire sur la partie en recouvrement~; ne pas vendre dans le mouvement, sauf si des limites l'y obligent~; réduire
les positions recouvrantes les moins liquides dès le début si cela s'impose~; envisager de fournir de la liquidité au vendeur si la part permanente est faible.
\end{solution}

\begin{solution}{iq:rs:capacity-decay-and-crowding:5}
Le profit est égal à la taille multipliée par le rendement net~: le rendement (et le ratio de Sharpe) baisse dès le premier dollar sous des coûts concaves, mais le profit croît tant que la taille croît plus vite
que le rendement baisse, jusqu'au point $A^* = (2g/3k)^2$.
\end{solution}

\begin{solution}{iq:rs:capacity-decay-and-crowding:6}
Si le vendeur échange $fQ_i$ de chaque titre sur une journée et que l'impact vaut $\eta\sigma_i\sqrt{fQ_i/V_i}$, un fonds détenant $h_i$ du titre perd
$\sum_i h_i\eta\sigma_i\sqrt{fQ_i/V_i} \propto \sqrt f$ sur la journée, dont la part permanente subsiste~; en étalant la vente sur $d$ jours et en laissant le reste se résorber
donne les trajectoires du chapitre.
\end{solution}
